\documentclass[preview]{standalone}

\usepackage{amsmath}
\usepackage{amssymb}
\usepackage{stellar}
\usepackage{definitions}

\begin{document}

\id{psicologia-significato-linguaggio}
\genpage

\section{Significato e linguaggio}

\begin{snippet}{significato-introduzione}
    Il significato è il fondamento della comunicazione. I significati non sono
    dati naturali immutabili: sono prodotti culturali e contingenti, generati
    dall'intreccio degli scambi tra parlanti.
\end{snippet}

\begin{snippetdefinition}{semantica-tratti-definition}{Semantica a tratti}
    La \emph{semantica a tratti} studia come i significati si compongono e si
    articolano attraverso insiemi di tratti interconnessi.
\end{snippetdefinition}

\begin{snippetdefinition}{semantica-vero-condizionale-definition}{Semantica vero-condizionale}
    La \emph{semantica vero-condizionale} sostiene che il significato di una
    parola o di una frase sia dato dal rapporto tra linguaggio e realtà, cioè
    dalle condizioni in base a cui un enunciato può essere considerato vero.
\end{snippetdefinition}

\begin{snippet}{semantica-vero-condizionale-limite}
    La concezione vero-condizionale presuppone che il significato consista in un
    insieme finito di condizioni necessarie e sufficienti. È stata criticata
    perché attribuisce al significato una stabilità oggettiva che non sempre
    corrisponde all'uso effettivo del linguaggio.
\end{snippet}

\begin{snippetdefinition}{semantica-strutturale-definition}{Semantica strutturale}
    La \emph{semantica strutturale} studia il linguaggio come rete autonoma di
    relazioni interne, in cui i fatti linguistici esistono l'uno in ragione
    dell'altro.
\end{snippetdefinition}

\begin{snippet}{valore-linguistico}
    Nella semantica strutturale, il significato dipende dal valore linguistico:
    una parola acquista significato attraverso le relazioni e le differenze che
    la separano dalle altre parole della lingua.
\end{snippet}

\begin{snippetdefinition}{semantica-cognitiva-definition}{Semantica cognitiva}
    La \emph{semantica cognitiva}, o interpretativa, sostiene che il significato
    consista nel modo in cui gli individui comprendono la propria esperienza.
\end{snippetdefinition}

\begin{snippet}{semantica-cognitiva-note}
    Per la semantica cognitiva, l'esperienza non coincide con la realtà in sé:
    nasce dall'incontro tra il punto di vista soggettivo del parlante e i dati
    esterni provenienti dal mondo.
\end{snippet}

\begin{snippetdefinition}{contesto-comunicativo-definition}{Contesto}
    Il \emph{contesto} è l'insieme delle restrizioni e delle opportunità
    biologiche, spaziotemporali, relazionali, istituzionali e culturali
    dell'\emph{hic et nunc} che, insieme a un dato testo, genera un messaggio
    dotato di senso.
\end{snippetdefinition}

\begin{snippet}{testo-contesto}
    La stabilità e la variabilità del significato dipendono
    dall'interdipendenza tra testo e contesto. Non vi è testo senza contesto e
    non vi è contesto senza testo: ciascuno influenza l'altro nella produzione
    del senso.
\end{snippet}

\begin{snippet}{linguaggio-universalismo-relativismo}
    Chomsky sostiene che il linguaggio abbia una struttura universale, innata e
    biologica. La prospettiva relativista, associata a Sapir e Whorf, sostiene
    invece che pensiero, lingua e cultura siano intrinsecamente intrecciati.
    Molti studiosi contemporanei adottano una versione debole del relativismo
    linguistico.
\end{snippet}

\begin{snippetdefinition}{fono-definition}{Fono}
    Un \emph{fono} è un suono che l'apparato vocale umano è in grado di emettere.
\end{snippetdefinition}

\begin{snippetdefinition}{fonema-definition}{Fonema}
    Il \emph{fonema} è un'unità fonica indivisibile e astratta, dotata di valore
    linguistico distintivo e oppositivo. Di per sé è privo di significato, ma
    permette di distinguere significati lessicali diversi.
\end{snippetdefinition}

\begin{snippetdefinition}{morfema-definition}{Morfema}
    Il \emph{morfema} è un'unità grammaticale minima, non ulteriormente
    scomponibile, isolabile e dotata di un proprio significato.
\end{snippetdefinition}

\begin{snippet}{morfemi-liberi-legati}
    I morfemi possono essere liberi, quando costituiscono parole autonome come
    \texttt{il} o \texttt{adesso}, oppure legati, quando richiedono la
    combinazione con altri morfemi, come \texttt{gatt-} e \texttt{-o} in
    \texttt{gatt-o}.
\end{snippet}

\begin{snippetdefinition}{lessico-definition}{Lessico}
    Il \emph{lessico} è l'insieme delle parole di una lingua.
\end{snippetdefinition}

\begin{snippetdefinition}{sintassi-definition}{Sintassi}
    La \emph{sintassi} è l'insieme organico e formale delle regole e dei
    procedimenti computazionali che consentono di disporre in ordine gerarchico
    gli elementi del lessico.
\end{snippetdefinition}

\begin{snippet}{creativita-linguistica}
    Ogni lingua possiede una grande creatività espressiva, perché permette di
    produrre un numero indefinito di enunciati. Tale creatività non è libera da
    vincoli, ma è governata da regole fonologiche, morfologiche, lessicali e
    sintattiche.
\end{snippet}

\end{document}
